%% ========================================================================
%% Chapter 9: Pemrograman Bash
%% ========================================================================

\chapter{Pemrograman Bash}
\label{ch:bash-programming}

Bab ini membahas dasar-dasar pemrograman Bash: membuat script, mengelola
variabel, menulis struktur kontrol, mendefinisikan fungsi, memproses argumen,
dan melakukan debugging. Konsep dibandingkan dengan Java agar mahasiswa yang
sudah mengenal pemrograman dapat membangun pemahaman lebih cepat.

\begin{examplebox}[Bash sebagai bahasa pemrograman]
Bayangkan Anda sudah pernah menulis program Java. Java pun memiliki
variabel, kondisi, perulangan, dan fungsi. Bash pun demikian, bedanya,
Bash menjalankan \textit{perintah sistem operasi} sebagai bagian dari logika
programnya. Sebuah script Bash adalah \textbf{daftar perintah terminal yang
disimpan dalam file}, dilengkapi logika kontrol agar dapat berjalan otomatis
tanpa keterlibatan manusia.
\end{examplebox}

%% ------------------------------------------------------------------------
\section{Dasar-dasar Scripting Bash}
\label{sec:membuat-script}

\subsection*{Cara Membuat dan Mengedit File Script}

Bab ini menggunakan dua cara membuat file: \textbf{nano} untuk menulis
interaktif, dan \textbf{heredoc} untuk membuat file langsung dari terminal.

\textbf{Menggunakan nano:}
\begin{enumerate}
  \item Buka atau buat file: \texttt{nano nama-file.sh}
  \item Ketik isi file kursor langsung aktif, tidak ada mode khusus.
  \item Tombol penting:
  \begin{itemize}
    \item \keystroke{Ctrl+O} lalu \keystroke{Enter} simpan file
    \item \keystroke{Ctrl+X} keluar dari nano
    \item \keystroke{Ctrl+W} cari teks dalam file
  \end{itemize}
\end{enumerate}

\begin{notebox}
Di bagian bawah layar nano, selalu tampil daftar shortcut aktif. Tanda
\texttt{\^{}} berarti \texttt{Ctrl}, jadi \texttt{\^{}O} = \keystroke{Ctrl+O}.
\end{notebox}

\textbf{Menggunakan heredoc} (untuk script pendek atau otomatisasi):
\begin{lstlisting}[language=bash, caption={Membuat file dengan heredoc}]
cat <<'EOF' > nama-file.sh
#!/bin/bash
echo "Halo dari script"
EOF
\end{lstlisting}

\subsection*{Struktur Dasar dan Shebang Line}

Setiap script Bash diawali dengan \textit{shebang line}:

\begin{lstlisting}[language=bash, caption={Struktur dasar script Bash}]
#!/bin/bash
# Komentar: jelaskan fungsi script ini

echo "Script berjalan"
\end{lstlisting}

Baris \texttt{\#!} memberi tahu kernel Linux interpreter mana yang
digunakan. Ada dua bentuk yang umum:

\begin{itemize}
  \item \texttt{\#!/bin/bash} path absolut ke Bash. Paling umum di
    Linux, sederhana, dan langsung.
  \item \texttt{\#!/usr/bin/env bash} meminta perintah \texttt{env}
    mencari Bash di \var{PATH}. Lebih portabel lintas sistem operasi.
\end{itemize}

\begin{center}
\begin{tabular}{lll}
\hline
\textbf{Aspek} & \textbf{\#!/bin/bash} & \textbf{\#!/usr/bin/env bash} \\
\hline
Lokasi Bash & Tetap \texttt{/bin/bash} & Dicari di \texttt{PATH} \\
Portabilitas & Linux standar & Linux dan macOS \\
Rekomendasi & Praktikum ini & Script lintas platform \\
\hline
\end{tabular}
\end{center}

\begin{notebox}
Shebang hanya bekerja jika file memiliki izin eksekusi (\texttt{chmod +x})
dan dijalankan langsung (\texttt{./script.sh}). Jika dipanggil
\texttt{bash script.sh}, shebang diabaikan karena interpreter
sudah ditentukan eksplisit.
\end{notebox}

Setelah membuat script, beri izin eksekusi dan jalankan:

\begin{lstlisting}[language=bash, caption={Memberi izin dan menjalankan script}]
chmod +x nama-file.sh
./nama-file.sh
\end{lstlisting}

\subsection*{Praktikum 7.1 Script Pertama: Laporan Sistem}

Tujuan: membuat, menyimpan, dan menjalankan script Bash pertama.

\begin{enumerate}
  \item Buat workspace praktikum:
\begin{lstlisting}[language=bash, caption={Menyiapkan workspace}]
mkdir -p ~/praktikum-os/week09/{scripts,logs,data}
cd ~/praktikum-os/week09/scripts
\end{lstlisting}

  \item Buat script dengan nano:
\begin{lstlisting}[language=bash, caption={Membuka nano}]
nano laporan-sistem.sh
\end{lstlisting}

  \item Ketik isi berikut, simpan (\keystroke{Ctrl+O} \keystroke{Enter}),
    lalu keluar (\keystroke{Ctrl+X}):
\begin{lstlisting}[language=bash, caption={Isi laporan-sistem.sh}]
#!/bin/bash
# Script: laporan-sistem.sh

echo "================================"
echo "  LAPORAN SISTEM"
echo "================================"
echo "Tanggal  : $(date '+%A, %d %B %Y')"
echo "Jam      : $(date '+%H:%M:%S')"
echo "Hostname : $(hostname)"
echo "User     : $(whoami)"
echo "CPU core : $(nproc)"
echo "RAM bebas: $(free -h | awk '/^Mem/ {print $4}')"
echo "Disk /   : $(df -h / | awk 'NR==2 {print $5}') terpakai"
echo "================================"
\end{lstlisting}

  \item Beri izin dan jalankan:
\begin{lstlisting}[language=bash, caption={Menjalankan script}]
chmod +x laporan-sistem.sh
./laporan-sistem.sh
\end{lstlisting}
\end{enumerate}

\begin{exercisebox}[Latihan 9.1]
Modifikasi \texttt{laporan-sistem.sh} agar menyimpan output ke file
\texttt{laporan-YYYY-MM-DD.txt} sekaligus menampilkannya di terminal.
Petunjuk: gunakan \texttt{tee} yang sudah dipelajari di bab sebelumnya.
\end{exercisebox}

%% ------------------------------------------------------------------------
\section{Variabel dan Parameter Posisional}
\label{sec:variabel-parameter}

\subsection*{Variabel di Bash vs Java}

Di Java: \texttt{String nama = "Budi";}. Di Bash, aturannya berbeda:

\begin{itemize}
  \item \textbf{Deklarasi}: \texttt{NAMA="Budi"} tanpa spasi di sekitar \texttt{=}
  \item \textbf{Akses nilai}: \texttt{\$NAMA} atau \texttt{\$\{NAMA\}}
  \item \textbf{Tipe data}: tidak ada. Semua variabel Bash adalah string
    secara default; aritmetika memerlukan sintaks khusus.
\end{itemize}

\begin{warningbox}
\texttt{NAMA = "Budi"} (dengan spasi) menghasilkan error Bash mengira
\texttt{NAMA} adalah nama perintah. Selalu tulis tanpa spasi.
\end{warningbox}

Variabel khusus yang paling sering dipakai dalam scripting:

\begin{center}
\begin{tabular}{ll}
\hline
\textbf{Variabel} & \textbf{Nilai} \\
\hline
\texttt{\$0} & Nama script yang sedang berjalan \\
\texttt{\$1}, \texttt{\$2}, \ldots & Argumen ke-1, ke-2, dst. \\
\texttt{\$\#} & Jumlah argumen yang diberikan \\
\texttt{\$@} & Semua argumen sebagai daftar terpisah \\
\texttt{\$?} & Exit code perintah terakhir (0 = sukses) \\
\texttt{\$\$} & PID script saat ini \\
\hline
\end{tabular}
\end{center}

\subsection*{Parameter Posisional dan Nilai Default}

Seperti metode Java yang menerima argumen
(\texttt{void fungsi(String a, String b)}), script Bash menerima argumen
dari command line. Argumen dibaca dengan \texttt{\$1}, \texttt{\$2}, dst.

Konstruk \texttt{\$\{1:-"default"\}} membaca \texttt{\$1}; jika kosong,
gunakan nilai default. Setara dengan
\texttt{args.length > 0 ? args[0] : "default"} di Java.

\subsection*{Aritmetika dan Substitusi Perintah}

Bash tidak langsung menghitung ekspresi aritmetika. Gunakan \texttt{\$(( ))}:

\begin{lstlisting}[language=bash, caption={Aritmetika dan substitusi perintah}]
A=10
B=3

echo "Jumlah  : $((A + B))"
echo "Modulo  : $((A % B))"
echo "Pembagian: $((A / B))"   # Hasil bulat, sama seperti int di Java

# Substitusi perintah: simpan output ke variabel
TANGGAL=$(date '+%F')
JUMLAH_FILE=$(ls | wc -l)
echo "Tanggal: $TANGGAL, File: $JUMLAH_FILE"
\end{lstlisting}

\begin{notebox}
Di Java, \texttt{10 / 3} untuk tipe \texttt{int} menghasilkan \texttt{3}.
Di Bash, \texttt{\$((10 / 3))} juga menghasilkan \texttt{3}. Gunakan
\texttt{bc} jika butuh presisi desimal.
\end{notebox}

\subsection*{Praktikum 7.2 Script Info Sistem dengan Argumen}

Tujuan: berlatih variabel, substitusi perintah, parameter posisional,
dan nilai default.

\begin{enumerate}
  \item Buat script:
\begin{lstlisting}[language=bash, caption={Membuat script info-sistem.sh}]
nano ~/praktikum-os/week09/scripts/info-sistem.sh
\end{lstlisting}

  \item Ketik isi berikut:
\begin{lstlisting}[language=bash, caption={Isi info-sistem.sh}]
#!/bin/bash
# Penggunaan: ./info-sistem.sh [nama-admin] [batas-disk-persen]

ADMIN=${1:-"Tidak dikenal"}
BATAS=${2:-80}
TANGGAL=$(date '+%F %T')
DISK_PERSEN=$(df / | awk 'NR==2 {print $5}' | tr -d '%')

echo "Admin    : $ADMIN"
echo "Tanggal  : $TANGGAL"
echo "Disk /   : ${DISK_PERSEN}% terpakai"
echo "Batas    : ${BATAS}%"

if [ "$DISK_PERSEN" -gt "$BATAS" ]; then
    echo "STATUS   : PERINGATAN - disk melebihi batas!"
else
    SISA=$((BATAS - DISK_PERSEN))
    echo "STATUS   : Normal (sisa toleransi ${SISA}%)"
fi
\end{lstlisting}

  \item Simpan, beri izin, uji dengan berbagai kombinasi argumen:
\begin{lstlisting}[language=bash, caption={Menguji script}]
chmod +x ~/praktikum-os/week09/scripts/info-sistem.sh
./info-sistem.sh
./info-sistem.sh "Dian" 50
./info-sistem.sh "Dian" 10
\end{lstlisting}
\end{enumerate}

\begin{exercisebox}[Latihan 9.2]
Buat script \texttt{kalkulator.sh} yang menerima tiga argumen:
\texttt{<angka1> <operator> <angka2>} dengan operator \texttt{+},
\texttt{-}, \texttt{*}, atau \texttt{/}. Contoh:
\texttt{./kalkulator.sh 20 + 5} menghasilkan \texttt{25}.
Gunakan \texttt{case} untuk memilih operasi, dan validasi jika argumen
tidak lengkap.
\end{exercisebox}

%% ------------------------------------------------------------------------
\section{Struktur Kontrol}
\label{sec:struktur-kontrol}

Struktur kontrol di Bash konsepnya sama dengan Java, hanya sintaksnya berbeda.

\begin{center}
\begin{tabular}{lll}
\hline
\textbf{Konsep} & \textbf{Java} & \textbf{Bash} \\
\hline
Kondisi if & \texttt{if (x > 5) \{} & \texttt{if [ \$x -gt 5 ]; then} \\
Blok kode & kurung kurawal & \texttt{then ... fi} \\
Else if & \texttt{else if (...) \{} & \texttt{elif ...; then} \\
Tidak sama & \texttt{!=} & \texttt{!=} (string), \texttt{-ne} (angka) \\
Dan / Atau & \texttt{\&\&} / \texttt{||} & \texttt{\&\&} / \texttt{||} \\
Switch & \texttt{switch (x) \{ case "a": \}} & \texttt{case \$x in "a") ... ;;} \\
\hline
\end{tabular}
\end{center}

\subsection*{Operator Perbandingan}

\textbf{Angka:} \texttt{-eq}, \texttt{-ne}, \texttt{-gt}, \texttt{-lt},
\texttt{-ge}, \texttt{-le}

\textbf{String:} \texttt{=} atau \texttt{==} (sama), \texttt{!=} (tidak sama),
\texttt{-z} (kosong), \texttt{-n} (tidak kosong)

\textbf{File:} \texttt{-f} (file ada), \texttt{-d} (direktori ada),
\texttt{-r} (bisa dibaca), \texttt{-x} (bisa dieksekusi)

\subsection*{Perulangan \texttt{for} dan \texttt{while}}

\begin{lstlisting}[language=bash, caption={Perulangan for dan while}]
# for iterasi atas daftar (mirip for-each Java)
for BULAN in Januari Februari Maret; do
    echo "Bulan: $BULAN"
done

# for C-style (mirip Java)
for ((I=1; I<=5; I++)); do
    echo "Angka: $I"
done

# for atas file
for FILE in ~/praktikum-os/week09/scripts/*.sh; do
    echo "Script: $(basename $FILE)"
done

# while selama kondisi benar
COUNTER=1
while [ $COUNTER -le 3 ]; do
    echo "Iterasi $COUNTER"
    COUNTER=$((COUNTER + 1))
done
\end{lstlisting}

\subsection*{Struktur \texttt{case}}

\texttt{case} setara \texttt{switch} di Java. Setiap blok diakhiri \texttt{;;}
(bukan \texttt{break}). Pola \texttt{*)} setara \texttt{default:}.

\begin{notebox}
Di Java, lupa menulis \texttt{break} menyebabkan \textit{fall-through}.
Di Bash, \texttt{;;} sudah otomatis menghentikan eksekusi blok
tidak ada fall-through secara default.
\end{notebox}

\subsection*{Praktikum 7.3 Script Grading dan Menu Interaktif}

Tujuan: menggabungkan \texttt{if}, \texttt{for}, \texttt{while}, dan
\texttt{case} dalam satu session.

\begin{enumerate}
  \item Buat script grading (menggunakan \texttt{if} dan \texttt{for}):
\begin{lstlisting}[language=bash, caption={Membuat script grading-batch.sh}]
nano ~/praktikum-os/week09/scripts/grading-batch.sh
\end{lstlisting}

  \item Ketik isi berikut:
\begin{lstlisting}[language=bash, caption={Isi grading-batch.sh}]
#!/bin/bash
# Script: grading-batch.sh
# Proses daftar nilai mahasiswa

MAHASISWA=("Andi:92" "Budi:73" "Citra:55" "Deni:80" "Eka:45")

echo "=== HASIL GRADING ==="
for ENTRI in "${MAHASISWA[@]}"; do
    NAMA=$(echo "$ENTRI" | cut -d: -f1)
    NILAI=$(echo "$ENTRI" | cut -d: -f2)

    if [ "$NILAI" -ge 85 ]; then
        GRADE="A"
    elif [ "$NILAI" -ge 75 ]; then
        GRADE="B"
    elif [ "$NILAI" -ge 65 ]; then
        GRADE="C"
    elif [ "$NILAI" -ge 55 ]; then
        GRADE="D"
    else
        GRADE="E"
    fi

    printf "%-10s %3d  %s\n" "$NAMA" "$NILAI" "$GRADE"
done
echo "====================="
\end{lstlisting}

  \item Simpan, beri izin, dan jalankan:
\begin{lstlisting}[language=bash, caption={Menjalankan script grading}]
chmod +x ~/praktikum-os/week09/scripts/grading-batch.sh
./grading-batch.sh
\end{lstlisting}

  \item Buat script menu interaktif (\texttt{while} + \texttt{case}):
\begin{lstlisting}[language=bash, caption={Membuat script menu-sistem.sh}]
nano ~/praktikum-os/week09/scripts/menu-sistem.sh
\end{lstlisting}

  \item Ketik isi berikut:
\begin{lstlisting}[language=bash, caption={Isi menu-sistem.sh}]
#!/bin/bash
# Menu interaktif pemantauan sistem

while true; do
    echo ""
    echo "===== MENU MONITOR ====="
    echo "1) Info disk"
    echo "2) Info memori"
    echo "3) Proses teratas"
    echo "4) Keluar"
    echo -n "Pilih [1-4]: "
    read PILIHAN

    case $PILIHAN in
        1) df -h ;;
        2) free -h ;;
        3) ps aux --sort=-%cpu | head -6 ;;
        4) echo "Sampai jumpa!"; exit 0 ;;
        *) echo "Pilihan tidak valid." ;;
    esac
done
\end{lstlisting}

  \item Beri izin dan jalankan, coba setiap opsi:
\begin{lstlisting}[language=bash, caption={Menjalankan menu}]
chmod +x ~/praktikum-os/week09/scripts/menu-sistem.sh
./menu-sistem.sh
\end{lstlisting}
\end{enumerate}

\begin{exercisebox}[Latihan 9.3]
Tambahkan ke script \texttt{grading-batch.sh} sebuah ringkasan di bagian
bawah yang menampilkan: jumlah mahasiswa per grade (A, B, C, D, E)
menggunakan perulangan \texttt{for} kedua yang mengiterasi array
\texttt{MAHASISWA}.
\end{exercisebox}

%% ------------------------------------------------------------------------
\section{Fungsi dalam Shell Script}
\label{sec:fungsi-shell-script}

\subsection*{Fungsi di Bash vs Java}

Di Java: \texttt{void namaMetode(String p1, String p2) \{ ... \}}.
Di Bash, fungsi didefinisikan tanpa daftar parameter eksplisit argumen
dikirim seperti command line dan dibaca dengan \texttt{\$1}, \texttt{\$2}, dst.

\begin{center}
\begin{tabular}{lll}
\hline
\textbf{Konsep} & \textbf{Java} & \textbf{Bash} \\
\hline
Definisi & \texttt{void f(String a, int b)} & \texttt{f() \{ ... \}} \\
Panggil & \texttt{f("x", 2)} & \texttt{f x 2} \\
Argumen & \texttt{a}, \texttt{b} & \texttt{\$1}, \texttt{\$2} \\
Return nilai & \texttt{return nilai} & \texttt{echo nilai} (via stdout) \\
Return status & exception & \texttt{return 0--255} \\
Variabel lokal & default lokal & wajib tulis \texttt{local} \\
\hline
\end{tabular}
\end{center}

\begin{notebox}
Di Bash, \texttt{return} hanya mengembalikan kode status (0--255). Untuk
mengembalikan nilai seperti string, gunakan \texttt{echo} di dalam fungsi
lalu tangkap dengan substitusi: \texttt{HASIL=\$(nama\_fungsi arg)}.
Di Java, variabel dalam metode otomatis lokal. Di Bash, tanpa \texttt{local},
variabel bocor ke lingkup global.
\end{notebox}

\begin{lstlisting}[language=bash, caption={Mendefinisikan dan memanggil fungsi}]
#!/bin/bash

# Fungsi mengembalikan nilai via stdout
hitung_luas() {
    local PANJANG=$1
    local LEBAR=$2
    echo $((PANJANG * LEBAR))
}

# Tangkap nilai dengan substitusi perintah
LUAS=$(hitung_luas 8 5)
echo "Luas: $LUAS cm2"

# Fungsi mengembalikan status sukses/gagal
bagi() {
    local A=$1 B=$2
    if [ "$B" -eq 0 ]; then
        echo "Error: pembagi nol" >&2
        return 1
    fi
    echo $((A / B))
}

HASIL=$(bagi 20 4) && echo "Hasil: $HASIL"
bagi 10 0 || echo "Exit code: $?"
\end{lstlisting}

\subsection*{Praktikum 7.4 Library Fungsi Validasi}

Tujuan: membuat file library fungsi yang dapat dimuat (\texttt{source})
oleh script lain konsep serupa dengan \textit{import} di Java.

\begin{enumerate}
  \item Buat file library:
\begin{lstlisting}[language=bash, caption={Membuat lib-validasi.sh}]
nano ~/praktikum-os/week09/scripts/lib-validasi.sh
\end{lstlisting}

  \item Ketik isi berikut:
\begin{lstlisting}[language=bash, caption={Isi lib-validasi.sh}]
#!/bin/bash
# lib-validasi.sh - Library fungsi validasi

adalah_angka() {
    [[ "$1" =~ ^[0-9]+$ ]]
}

file_bisa_dibaca() {
    [ -f "$1" ] && [ -r "$1" ]
}

error_exit() {
    echo "ERROR: $1" >&2
    exit 1
}

info()   { echo "[INFO] $1"; }
sukses() { echo "[OK]   $1"; }
\end{lstlisting}

  \item Buat script yang menggunakan library:
\begin{lstlisting}[language=bash, caption={Membuat script pakai-library.sh}]
nano ~/praktikum-os/week09/scripts/pakai-library.sh
\end{lstlisting}

  \item Ketik isi berikut:
\begin{lstlisting}[language=bash, caption={Isi pakai-library.sh}]
#!/bin/bash
# Muat library (seperti import di Java)
source ~/praktikum-os/week09/scripts/lib-validasi.sh

ANGKA=$1
FILE=$2

[ -z "$ANGKA" ] || [ -z "$FILE" ] && \
    error_exit "Penggunaan: $0 <angka> <path-file>"

if adalah_angka "$ANGKA"; then
    sukses "Input '$ANGKA' adalah angka valid"
else
    error_exit "'$ANGKA' bukan angka"
fi

if file_bisa_dibaca "$FILE"; then
    sukses "File '$FILE' bisa dibaca"
    info "Jumlah baris: $(wc -l < "$FILE")"
else
    error_exit "File '$FILE' tidak ada atau tidak bisa dibaca"
fi
\end{lstlisting}

  \item Beri izin dan uji semua skenario:
\begin{lstlisting}[language=bash, caption={Menguji semua skenario}]
chmod +x ~/praktikum-os/week09/scripts/pakai-library.sh
./pakai-library.sh 42 /etc/hostname
./pakai-library.sh abc /etc/hostname
./pakai-library.sh 42 /tidak-ada.txt
./pakai-library.sh
\end{lstlisting}
\end{enumerate}

\begin{exercisebox}[Latihan 9.4]
Tambahkan fungsi \texttt{konfirmasi()} ke \texttt{lib-validasi.sh}.
Fungsi ini menampilkan pertanyaan, membaca input Y/N dari user,
mengembalikan \texttt{0} jika Y dan \texttt{1} jika N. Buat script demo
yang memanggil fungsi ini sebelum menghapus sebuah file.
\end{exercisebox}

%% ------------------------------------------------------------------------
\section{Pengolahan Argumen Command Line}
\label{sec:pengolahan-argumen}

Script yang baik memproses argumen secara fleksibel: validasi jumlah
argumen, dan opsi seperti \texttt{-v} atau \texttt{-o file}.

\texttt{getopts} adalah cara standar Bash memproses opsi mirip cara
kerja perintah Unix seperti \texttt{ls -la} atau \texttt{grep -i}:

\begin{lstlisting}[language=bash, caption={Pola dasar getopts}]
#!/bin/bash
VERBOSE=false
OUTPUT_FILE=""

while getopts "vo:" OPSI; do
    case $OPSI in
        v) VERBOSE=true ;;
        o) OUTPUT_FILE="$OPTARG" ;;
        *) echo "Penggunaan: $0 [-v] [-o file]"; exit 1 ;;
    esac
done
shift $((OPTIND - 1))   # Geser argumen agar $1 = argumen non-opsi pertama

echo "Verbose : $VERBOSE"
echo "Output  : ${OUTPUT_FILE:-tidak ada}"
echo "Sisa    : $@"
\end{lstlisting}

Dalam \texttt{"vo:"}, tanda \texttt{:} setelah huruf berarti opsi itu
memerlukan nilai (\texttt{-o nama-file}), bukan sekadar flag (\texttt{-v}).

\subsection*{Praktikum 7.5 Script Backup dengan Opsi}

Tujuan: membuat script yang memproses opsi \texttt{getopts} dan
mengintegrasikan semua konsep sebelumnya.

\begin{enumerate}
  \item Buat script:
\begin{lstlisting}[language=bash, caption={Membuat backup-data.sh}]
nano ~/praktikum-os/week09/scripts/backup-data.sh
\end{lstlisting}

  \item Ketik isi berikut:
\begin{lstlisting}[language=bash, caption={Isi backup-data.sh}]
#!/bin/bash
# Penggunaan: ./backup-data.sh [-v] [-c] [-l logfile] <sumber> <tujuan>

VERBOSE=false
COMPRESS=false
LOG_FILE=""

while getopts "vcl:" OPSI; do
    case $OPSI in
        v) VERBOSE=true ;;
        c) COMPRESS=true ;;
        l) LOG_FILE="$OPTARG" ;;
        *) echo "Penggunaan: $0 [-v] [-c] [-l logfile] <sumber> <tujuan>"
           exit 1 ;;
    esac
done
shift $((OPTIND - 1))

SUMBER=$1
TUJUAN=$2

log() {
    local MSG="[$(date '+%T')] $1"
    echo "$MSG"
    [ -n "$LOG_FILE" ] && echo "$MSG" >> "$LOG_FILE"
}

[ -z "$SUMBER" ] || [ -z "$TUJUAN" ] && {
    echo "ERROR: sumber dan tujuan wajib diisi"; exit 1; }

[ ! -d "$SUMBER" ] && { log "ERROR: $SUMBER tidak ada"; exit 2; }

mkdir -p "$TUJUAN"
TANGGAL=$(date '+%F-%H%M%S')
[ "$VERBOSE" = true ] && log "Sumber: $SUMBER | Tujuan: $TUJUAN"

if [ "$COMPRESS" = true ]; then
    ARSIP="$TUJUAN/backup-$(basename "$SUMBER")-$TANGGAL.tar.gz"
    tar -czf "$ARSIP" -C "$(dirname "$SUMBER")" "$(basename "$SUMBER")"
    log "Arsip: $ARSIP ($(du -sh "$ARSIP" | cut -f1))"
else
    cp -r "$SUMBER" "$TUJUAN/backup-$(basename "$SUMBER")-$TANGGAL"
    log "Backup selesai."
fi
\end{lstlisting}

  \item Beri izin dan uji:
\begin{lstlisting}[language=bash, caption={Menguji script backup}]
chmod +x ~/praktikum-os/week09/scripts/backup-data.sh
cd ~/praktikum-os/week09/scripts

# Tanpa opsi
./backup-data.sh ~/praktikum-os/week09/data ~/praktikum-os/week09/logs

# Dengan verbose dan kompresi + log ke file
./backup-data.sh -v -c -l ../logs/backup.log \
  ~/praktikum-os/week09/data ~/praktikum-os/week09/logs

cat ../logs/backup.log
\end{lstlisting}
\end{enumerate}

%% ------------------------------------------------------------------------
\section{Debugging Shell Script}
\label{sec:debugging-script}

Debugging Bash berbeda dari Java yang memiliki debugger interaktif.
Bash menyediakan mekanisme tracing untuk memperlihatkan setiap baris yang
dieksekusi.

\begin{center}
\begin{tabular}{lll}
\hline
\textbf{Flag} & \textbf{Arti} & \textbf{Cara aktifkan} \\
\hline
\texttt{-x} & Tampilkan setiap perintah sebelum dijalankan & \texttt{bash -x script.sh} \\
\texttt{-e} & Keluar jika ada perintah gagal & \texttt{set -e} \\
\texttt{-u} & Keluar jika variabel tidak terdefinisi & \texttt{set -u} \\
\texttt{-n} & Cek sintaks tanpa menjalankan & \texttt{bash -n script.sh} \\
\hline
\end{tabular}
\end{center}

Kombinasi \texttt{set -euo pipefail} di awal script menangkap banyak
kesalahan umum secara otomatis. Tambahkan \texttt{trap} untuk cleanup:

\begin{lstlisting}[language=bash, caption={Pola script defensif}]
#!/bin/bash
set -euo pipefail

cleanup() { echo "[CLEANUP] selesai."; }
trap cleanup EXIT      # cleanup dipanggil selalu saat script berakhir

DEBUG=${DEBUG:-false}
debug() { [ "$DEBUG" = true ] && echo "[DEBUG] $*" >&2; }

debug "Script dimulai dengan PID $$"
echo "Script berjalan normal"
\end{lstlisting}

\begin{notebox}
\texttt{trap cleanup EXIT} setara dengan \texttt{finally} di Java
memastikan cleanup dijalankan meskipun script gagal di tengah jalan.
\end{notebox}

\subsection*{Praktikum 7.6 Debugging Script}

Tujuan: menggunakan teknik debugging untuk menganalisis dan memperbaiki
script.

\begin{enumerate}
  \item Buat script untuk dianalisis:
\begin{lstlisting}[language=bash, caption={Membuat debug-latihan.sh}]
nano ~/praktikum-os/week09/scripts/debug-latihan.sh
\end{lstlisting}

  \item Ketik isi berikut:
\begin{lstlisting}[language=bash, caption={Isi debug-latihan.sh}]
#!/bin/bash
# Script: debug-latihan.sh
# Penggunaan: ./debug-latihan.sh <direktori> <batas-MB>

DIREKTORI=$1
BATAS=$2

if [ $# -ne 2 ]; then
    echo "Penggunaan: $0 <direktori> <batas-MB>"
    exit 1
fi

UKURAN=$(du -sm "$DIREKTORI" | cut -f1)

echo "Direktori : $DIREKTORI"
echo "Ukuran    : ${UKURAN} MB"
echo "Batas     : ${BATAS} MB"

if [ "$UKURAN" -gt "$BATAS" ]; then
    echo "PERINGATAN: Ukuran melebihi batas!"
    echo "Kelebihan: $((UKURAN - BATAS)) MB"
else
    echo "Status: Normal (sisa: $((BATAS - UKURAN)) MB)"
fi
\end{lstlisting}

  \item Cek sintaks, lalu jalankan dengan tracing:
\begin{lstlisting}[language=bash, caption={Menganalisis script dengan debugging}]
chmod +x ~/praktikum-os/week09/scripts/debug-latihan.sh
bash -n debug-latihan.sh && echo "Sintaks OK"
bash -x debug-latihan.sh /etc 10
./debug-latihan.sh /var 50
./debug-latihan.sh
\end{lstlisting}
\end{enumerate}

\begin{exercisebox}[Latihan 9.5]
Script \texttt{debug-latihan.sh} tidak menangani direktori yang tidak ada.
Perbaiki dengan menambahkan:
\begin{itemize}
  \item \texttt{set -e} di baris kedua
  \item Pengecekan \texttt{-d "\$DIREKTORI"} sebelum memanggil \texttt{du}
  \item Pesan error yang informatif jika direktori tidak ditemukan
\end{itemize}
Uji dengan direktori yang tidak ada.
\end{exercisebox}

%% ------------------------------------------------------------------------
\section{Best Practices}
\label{sec:best-practices}

Script Bash yang baik mengikuti prinsip berikut:

\begin{enumerate}
  \item \textbf{Selalu tulis shebang} (\texttt{\#!/bin/bash}).
  \item \textbf{Gunakan \texttt{set -euo pipefail}} di awal script.
  \item \textbf{Kutip semua variabel}: \texttt{"\$NAMA"}, bukan \texttt{\$NAMA}
    mencegah masalah jika nilai mengandung spasi.
  \item \textbf{Gunakan \texttt{local}} di dalam fungsi untuk mencegah
    variabel bocor ke global.
  \item \textbf{Validasi semua input} sebelum diproses.
  \item \textbf{Nama variabel deskriptif}: \texttt{DIREKTORI\_BACKUP}
    lebih jelas dari \texttt{D}.
  \item \textbf{Gunakan \texttt{trap}} untuk cleanup saat script berakhir.
  \item \textbf{Pisahkan konfigurasi dari logika}: taruh variabel konfigurasi
    di bagian atas script.
\end{enumerate}

Template script profesional yang menerapkan semua prinsip di atas:

\begin{longlisting}[language=bash, caption={Template script Bash profesional}]
#!/bin/bash
# ====================================================
# Script: nama-script.sh
# Deskripsi: Penjelasan singkat fungsi script ini
# Penggunaan: ./nama-script.sh [-v] [-h] <argumen>
# ====================================================

set -euo pipefail

SCRIPT_NAME=$(basename "$0")
VERBOSE=false

usage() {
    echo "Penggunaan: $SCRIPT_NAME [-v] [-h] <argumen>"
    echo "  -v  Verbose mode"
    echo "  -h  Tampilkan bantuan"
    exit 0
}

log()         { echo "[$SCRIPT_NAME] $*"; }
log_verbose() { [ "$VERBOSE" = true ] && log "$*"; }
error_exit()  { echo "ERROR: $*" >&2; exit 1; }
cleanup()     { log_verbose "Cleanup selesai."; }

trap cleanup EXIT

while getopts "vh" OPSI; do
    case $OPSI in
        v) VERBOSE=true ;;
        h) usage ;;
        *) error_exit "Opsi tidak dikenal. Gunakan -h." ;;
    esac
done
shift $((OPTIND - 1))

[ $# -lt 1 ] && error_exit "Argumen tidak lengkap. Gunakan -h."

log "Script dimulai"
log_verbose "Argumen: $*"

# Logika utama di sini

log "Selesai."
\end{longlisting}

%% ------------------------------------------------------------------------
\section{Tugas Praktikum}
\label{sec:tugas-praktikum-bash-prog}

\begin{warningbox}
Kerjakan tugas di \texttt{\$HOME/praktikum-os/week09/}. Script di
\texttt{scripts/}, output di \texttt{logs/}.
\end{warningbox}

\subsection*{Tugas 1 Script Absensi Kelas}

\textbf{Konteks:} instruktur mencatat kehadiran mahasiswa dari command line.

\textbf{Instruksi:}
\begin{enumerate}
  \item Buat script \texttt{absensi.sh} yang:
  \begin{itemize}
    \item Menerima argumen nama mahasiswa dan status
      (\texttt{hadir}/\texttt{izin}/\texttt{alpha})
    \item Menyimpan entri ke \texttt{absensi-YYYY-MM-DD.txt}
      dengan format \texttt{[HH:MM] NAMA - STATUS}
    \item Opsi \texttt{-r}: tampilkan rekapitulasi (jumlah per status)
    \item Opsi \texttt{-h}: tampilkan bantuan
  \end{itemize}
  \item Rekam minimal 5 entri dan tampilkan rekapitulasinya.
\end{enumerate}

\textbf{Konsep wajib:} variabel, parameter posisional, \texttt{getopts},
\texttt{if}, \texttt{for}, fungsi, dan redirection ke file.

\subsection*{Tugas 2 Script Health Check Sistem}

\textbf{Konteks:} administrator membuat pemeriksaan kondisi server
sebelum maintenance.

\textbf{Instruksi:}
\begin{enumerate}
  \item Buat script \texttt{healthcheck.sh} menggunakan template
    profesional dari bagian Best Practices.
  \item Script menampilkan: tanggal/waktu, hostname, uptime, penggunaan
    CPU, memori, dan disk untuk setiap filesystem yang terpasang.
  \item Jika penggunaan disk mana pun melebihi 80\%, tampilkan peringatan.
  \item Simpan hasil ke \texttt{healthcheck-YYYY-MM-DD.log} dan tampilkan
    ke terminal sekaligus menggunakan \texttt{tee}.
  \item Opsi \texttt{-t <persen>} mengubah batas peringatan disk
    (default 80).
\end{enumerate}

\textbf{Konsep wajib:} \texttt{set -euo pipefail}, \texttt{trap},
\texttt{getopts}, fungsi dengan \texttt{local}, \texttt{for},
\texttt{if}, dan \texttt{tee}.

%% ------------------------------------------------------------------------
\section{Rangkuman}
\label{sec:rangkuman-bash-prog}

\begin{itemize}
  \item Shebang (\texttt{\#!/bin/bash}) menentukan interpreter. Tanpa
    shebang, perilaku script tidak terdefinisi.
  \item Variabel Bash tidak bertipe; aritmetika pakai \texttt{\$(( ))};
    output perintah ditangkap dengan \texttt{\$( )}.
  \item Parameter posisional (\texttt{\$1}, \texttt{\$\#}, \texttt{\$@})
    memungkinkan script menerima argumen dari command line.
  \item Struktur kontrol \texttt{if}, \texttt{for}, \texttt{while}, dan
    \texttt{case} setara dengan Java, hanya sintaksnya berbeda.
  \item Fungsi menerima argumen via \texttt{\$1}/\texttt{\$2}; gunakan
    \texttt{local} untuk variabel lokal dan \texttt{echo} untuk
    mengembalikan nilai.
  \item \texttt{getopts} memproses opsi \texttt{-x} seperti perintah Unix
    standar.
  \item \texttt{set -euo pipefail} dan \texttt{trap} adalah fondasi script
    yang andal dan mudah di-debug.
\end{itemize}

%% ------------------------------------------------------------------------
\section{Referensi}
\label{sec:referensi-bash-prog}

\begin{enumerate}
  \item GNU Bash Reference Manual \url{https://www.gnu.org/software/bash/manual/}
  \item William Shotts, \textit{The Linux Command Line}, 2nd Edition. No Starch Press.
  \item Arnold Robbins dan Nelson H.F.\ Beebe, \textit{Classic Shell Scripting}. O'Reilly.
  \item Richard Blum dan Christine Bresnahan, \textit{Linux Command Line and Shell Scripting Bible}, 4th Edition.
  \item ShellCheck alat analisis statis untuk script Bash: \url{https://www.shellcheck.net}
\end{enumerate}
